\documentclass[10pt,a4paper]{article}
\usepackage{fontspec,amsmath,amssymb,graphicx}
\usepackage[margin=18mm]{geometry}
\setmainfont{Times New Roman}
\setlength{\parindent}{0pt}
\setlength{\parskip}{5pt}
\allowdisplaybreaks
\sloppy
\emergencystretch=2em
\begin{document}
\section*{BT\_5\_5\_3\_13 — Lời giải bổ sung}
(a) Phát biểu thiếu giả thiết $\Delta u=0$: với $u$ chỉ $C^2$ tùy ý, công thức đúng là $\int_{\partial\Omega}u\partial_\nu u=\int_\Omega(|\nabla u|^2+u\Delta u)$, không luôn không âm. Phản ví dụ trên quả cầu đơn vị: $u=2-|x|^2$, tích phân biên $=-2|\partial B|<0$. Nếu thêm điều hòa, tích phân bằng $\int|\nabla u|^2\ge0$, bằng 0 đúng khi $u$ hằng.
(b) Hiệu $w$ của hai nghiệm cùng $f,\beta$ là điều hòa và $\partial_\nu w=-\alpha w$. Vì $\alpha>0$ hầu khắp, $\int|\nabla w|^2=-\int\alpha w^2\le0$ suy $w=0$. Vậy bài Robin có nhiều nhất một nghiệm.\end{document}